% Original, logo-free starter based on the CC BY Tsinghua letterhead example.
\documentclass[12pt]{article}
\usepackage{fontspec}
\usepackage{xeCJK}
\usepackage[left=3cm,right=3cm,top=4.2cm,bottom=3cm,footskip=2cm,headheight=2.2cm,headsep=1cm]{geometry}
\usepackage{fancyhdr}
\usepackage{tabularx}
\usepackage[hidelinks]{hyperref}
\usepackage{zhnumber}
\setmainfont{TeX Gyre Termes}
\setCJKmainfont{FandolSong-Regular.otf}

\fancypagestyle{plain}{%
  \fancyhf{}
  \fancyhead[L]{\sffamily\small Tsinghua University}
  \fancyhead[R]{\sffamily\small Department of Mathematical Sciences}
  \renewcommand{\footrulewidth}{.4pt}
  \fancyfoot[L]{\footnotesize
    Prof. Chen Wei \\ Department of Mathematical Sciences \\ Beijing 100084, China \\
    Email: \href{mailto:chen.wei@example.edu}{chen.wei@example.edu}}
}
\pagestyle{plain}
\pagenumbering{gobble}

\begin{document}
\begin{flushright}
  Beijing, \zhtoday
\end{flushright}

\vspace{1.2cm}
{\Large\noindent 尊敬的李教授：}

\vspace{1em}
\large
\setlength{\parindent}{2em}

感谢您邀请我们参加数学教育工作组的秋季研讨会。我们很高兴分享近期
关于开放教材和可复现实验的工作，并与贵团队讨论如何为研究生课程设计
一套可持续的资料维护流程。

我们建议安排一次四十五分钟的线上会议，重点讨论课程目标、资料版本
管理和学生反馈。随信附上简短的议程草案；如时间合适，我们将在会议前
发送示例章节和实验记录。

感谢您的考虑，期待在九月继续交流。

\vspace{2em}
\begin{flushright}
  Chen Wei\\
  Department of Mathematical Sciences\\
  Tsinghua University
\end{flushright}
\end{document}
